\section{Jensen--Shannon 散度与尺度选择}
\label{sec:07-Information-Theory/03-Divergences/03}

你有两个新闻语料库，想比较它们的用词分布。某个专有名词——比如一家只在第一个语料库所在国家运营的公司名——在第一个语料库里出现了几千次，在第二个语料库里一次都没出现过。第二个语料库在该词上的概率是零。

KL 散度直接就不工作了。任一方向都可能因为支持不重合而变成无穷。而且``两个语料库有多不同''是一个没有天然方向的问题——你不知道该把哪个库当作参考。

\subsection{先试试笨办法：把两个方向的 KL 平均}

最直接的修补是把两个方向的 KL 取平均：

\[
\frac12 \KL(P \| Q) + \frac12 \KL(Q \| P).
\]

它确实对称了。但支持不重合的问题还在——只要存在某个词只在 \(P\) 中出现而 \(Q\) 中不出现，\(\KL(P \| Q)\) 的对应项就是无穷，平均值跟着无穷。双向平均没有解决根本问题，只是把爆炸从单侧变成了两侧。

真正的问题出在：两个分布直接把对方当作参考分布。在参考分布为零的地方，被比较的分布必须也为零，否则 KL 发散。如果两者的支持本来就不完全重叠，这种``互相做参考''的方式注定要炸。

\subsection{换一个思路：先混合，再分别比较}

Jensen--Shannon 散度的想法很朴素：不要拿 \(P\) 和 \(Q\) 直接互相对照。先构造一个双方都能覆盖的共同参照，再各自与这个参照比较。

令

\[
M = \frac12 (P + Q).
\]

\(M\) 是 \(P\) 和 \(Q\) 的等权混合。因为 \(M \ge P/2\) 且 \(M \ge Q/2\)，在相应支持上有密度比 \(dP/dM \le 2\) 与 \(dQ/dM \le 2\)。两个 KL 项都不超过 \(\log 2\)——无论 \(P\) 和 \(Q\) 的支持差多远，跟 \(M\) 比较都不会发散。这是在连续分布上也成立的论证，不依赖逐点概率。

Jensen--Shannon 散度定义为

\[
\operatorname{JS}(P, Q) = \frac12 \KL(P \| M) + \frac12 \KL(Q \| M).
\]

定义直接送出三个 KL 没有同时具备的性质：

\begin{itemize}
\tightlist
\item
  对称：交换 \(P, Q\) 不改变 \(M\)，结果不变；
\item
  非负：两个 KL 项都非负；
\item
  有界：以 2 为对数底时，\(0 \le \operatorname{JS}(P, Q) \le 1\) bit。
\end{itemize}

左端等号在 \(P = Q\) 时取得——两分布相同，各自与混合的 KL 都是零。右端在 \(P\) 和 \(Q\) 支持完全不相交时达到——此时观察一个样本就足以判断它来自哪个分布。

\subsection{1 bit 的上界不是凑出来的}

0 到 1 bit 这个范围，背后有一个清楚的解释。引入一个公平的来源标签 \(Z \in \{0, 1\}\)：

\[
X \mid Z = 0 \sim P, \qquad X \mid Z = 1 \sim Q.
\]

先抽来源（公平硬币），再从对应的分布里抽 \(X\)。如果忽略来源，\(X\) 的边缘分布正是 \(M\)。现在看互信息 \(I(Z; X)\)——看到样本 \(X\) 后，你对来源 \(Z\) 的不确定性减少了多少。

展开 \(I(Z; X) = \KL(P_{ZX} \| P_Z P_X)\)。联合概率拆成 \(P(z) P(x \mid z)\)，公共因子 \(P(z)\) 在对数里消去：

\[
\begin{aligned}
I(Z; X)
&= \sum_z P(z) \; \KL(P_{X\mid Z=z} \| P_X) \\
&= \frac12 \KL(P \| M) + \frac12 \KL(Q \| M) \\
&= \operatorname{JS}(P, Q).
\end{aligned}
\]

因此 JS 散度的含义非常直白：它就是看到样本后，你获得了多少关于``它来自 \(P\) 还是 \(Q\)''的信息。来源标签本身只有 \(H(Z) = 1\) bit 的熵——一个公平硬币的不确定性最多是 1 bit。互信息不可能超过边际熵，所以

\[
\operatorname{JS}(P, Q) = I(Z; X) \le H(Z) = 1 \text{ bit}.
\]

上界不是技巧的产物，是来源标签只携带一 bit 信息这个简单事实的后果。

\subsection{两个极端例子把尺度锚定}

第一个极端：\(P = Q\)。此时 \(M = P = Q\)，两个 KL 都为零，\(\operatorname{JS} = 0\)。样本对判断来源毫无帮助——两个分布完全相同。

第二个极端：

\[
P = (1, 0), \qquad Q = (0, 1).
\]

双向 KL 都是无穷，但混合 \(M = (1/2, 1/2)\) 让一切有限：

\[
\KL(P \| M) = 1 \cdot \log_2 \frac{1}{1/2} + 0 \cdot \log_2 0 = 1 \text{ bit},
\]
\[
\KL(Q \| M) = 1 \text{ bit},
\]
\[
\operatorname{JS}(P, Q) = 1 \text{ bit}.
\]

其中 \(\KL(P \| M)\) 正是第一节用 \(P = (1, 0)\) 与 \(M = (1/2, 1/2)\) 算出的那个 1 bit。若 \(P\) 和 \(Q\) 的支持完全不交，看到样本就已经完全暴露来源标签，\(Z\) 的 1 bit 不确定性被彻底消除。

\subsection{非等权混合：来源先验不必公平}

如果来源先验不是一半一半——比如第一个语料库的样本量是第二个的两倍——令 \(P(Z = 0) = \pi\)，混合分布为

\[
M_\pi = \pi P + (1 - \pi) Q.
\]

相应的加权 JS 是

\[
\operatorname{JS}_\pi(P, Q) = \pi \KL(P \| M_\pi) + (1 - \pi) \KL(Q \| M_\pi) = I(Z; X).
\]

上界不再是固定的 1 bit，而是来源标签的二元熵 \(h_2(\pi)\)——先验本身的不确定性决定了互信息的上限。如果 \(\pi = 0.9\)（来源极不均衡），\(h_2(0.9) \approx 0.469\) bit，即使两个分布完全不同，互信息也不会超过这个值——样本大概率来自主导分布，观察对判断来源的帮助比均衡情形小。

这说明散度的数值并非绝对：同样的分布对，配上不同的混合权重、对数底和采样机制，给出的数字和上界都不同。报告 JS 值时，至少要说清权重和底。

\subsection{JS 也不是万能距离}

JS 在对称性和有界性上比 KL 友好，但它一般仍不满足三角不等式，严格说是散度而非度量。不过 \(\sqrt{\operatorname{JS}}\) 是一个度量——这是一个已证明的定理，此处只引用结论。

更广泛的族是 \(f\)-散度，它统一了 KL、JS、总变差距离、Hellinger 距离等。不同散度站在不同视角观察``分布差多少''，不存在一个对所有任务都最优的通用散度。

选择比较量前，可以依次追问几个问题：

\begin{enumerate}
\tightlist
\item
  是否存在明确的真实分布与近似分布？若有，KL 的方向可能正是问题的一部分——\(P\) 是真实数据，\(Q\) 是模型近似，这时 \(\KL(P \| Q)\) 的方向有清晰的操作含义。
\item
  支持是否可能不重合？若会，直接 KL 可能成为无穷，JS 或其它有界散度更安全。
\item
  是否需要对称性、有限上界或三角不等式？聚类和可视化通常需要对称性；需要几何解释时可能要求度量性质。
\item
  你关心的是编码冗余（KL）、分类可辨性（JS）、概率质量差（总变差），还是样本空间中的搬运成本（Wasserstein）？
\end{enumerate}

例如：比较概率预测与真实生成机制时，\(\KL(P \| Q)\) 直接就是负对数损失的期望超过量。比较两个地位相同的数据集词分布时，JS 的方向无关性更自然。若两个图像分布只发生了轻微空间平移，Wasserstein 距离能捕捉到``近但未重叠''的几何邻近性，KL 和 JS 却可能饱和到接近各自的上界。

脱离任务说``某个散度最好''，通常没有意义。尺子本身不带方向——你需要先知道要量什么。

\subsection{实际中，散度是算不精确的}

全部公式都假设 \(P\) 和 \(Q\) 已知。但数据分析里通常只有有限样本。直接把频率代入公式：

\begin{itemize}
\tightlist
\item
  许多词或值在样本中未出现，产生大量零概率，KL 对此极其敏感；
\item
  随意加平滑（Laplace 平滑、Good--Turing 等）会改变目标分布，从估计散度变成估计被平滑扭曲后的量；
\item
  高维连续数据中，密度估计本身的误差可能远超散度计算精度。
\end{itemize}

因此报告经验散度时，至少需要交代：
\begin{itemize}
\tightlist
\item
  分布是如何得到的——已知模型、直方图、核密度估计还是频率代入；
\item
  零概率和未观察类别如何处理；
\item
  使用哪个方向、对数底、混合权重；
\item
  数值是否带有抽样不确定性（如 bootstrap 区间或多次划分的波动范围）。
\end{itemize}

数学定义给了比较尺度，但不会替任何人解决统计估计问题。这个边界在后面的数理统计部分会再次出现——届时我们会看到，有些看起来``算出来''的散度值，实际可能是估计方法的假象。
